\section{Practical guidance}

\begin{enumerate}[leftmargin=1.6em]
\item \textbf{Read importance ordinally, not as an effect size.} Treat it as a
  ranked shortlist among comparable features, never as a magnitude to line up
  against OLS coefficients. A high bar may reflect a marginal effect, an
  interaction (Section~\ref{sec:interactions}), a non-linear effect
  (Section~\ref{sec:nonlinear}), or merely a skewed feature distribution
  (Section~\ref{sec:confound}).
\item \textbf{Disambiguate with partial-dependence plots.} PDPs \citep{friedman2001}
  recover exactly what importance discards: in Section~\ref{sec:confound} the skewed-feature
  model's PDPs are straight lines while the non-linear model's curve --- same
  importances, different shapes. For the Section~\ref{sec:interactions}
  interaction, the one-way PDPs over $a$ and $b$ are flat --- the parity pattern
  only appears in the joint $(a,b)$ view --- which is itself the tell that the
  effect is a higher-order interaction that per-feature importance has flattened.
  Curvature (or a suspiciously flat low-order PDP over a high-importance variable)
  is a concrete cue to add a polynomial or interaction term back into an
  interpretable model.

  \emph{With one caveat.} A PDP averages the model over the marginal
  distribution of the other features, which is exact here because our features
  are independent. When features are dependent it evaluates the model on
  combinations that do not occur in the data, and \citet{hooker2021} show that
  the resulting curve can report the model's
  extrapolation rather than any effect --- the same mechanism that inflates
  permutation importance for correlated predictors. On correlated data, use
  ALE \citep{apley2020}, which accumulates local effects conditional on the
  neighbouring feature distribution instead of marginalising across it.
\item \textbf{\code{impurity\_corrected} fixes a different problem.} It removes
  the bias toward null / high-cardinality features (use it when feature types
  are mixed), and it tracks the variance contribution more faithfully than plain
  impurity, which the default \code{mtry} compresses (Section~\ref{sec:linear})
  --- at the price of far higher seed-to-seed variance on weak features, so
  read its ratios only as averages over seeds.
  What it does not change is the \emph{units}: it is still a variance, not a
  per-unit effect, so nothing in
  Sections~\ref{sec:interactions}--\ref{sec:confound} goes away --- the table
  in Section~\ref{sec:confound} shows it pinned across the two worlds exactly,
  like plain impurity. Debiased
  MDI variants exist \citep{li2019,zhou2021} but target the same type-bias, not
  the shape/scale reading examined here.
\item \textbf{Switching to permutation importance does not fix the shape/skew
  conflation.} Permutation and conditional \citep{strobl2008} importances address
  other problems (and are worth using), but as Section~\ref{sec:confound} shows
  they too summarise variance contribution and cannot separate effect shape from
  feature distribution. Only a shape-aware view (PDP/ALE \citep{friedman2001,apley2020},
  SHAP \citep{lundberg2017}) does that.
\item \textbf{If you know the structure is linear, use the linear model.} Random
  forests buy you interaction- and non-linearity-detection; when there is nothing
  to detect, OLS is both more accurate and directly interpretable.
\end{enumerate}
